\documentclass[10pt,a4paper]{amsart}
\usepackage[margin=19mm]{geometry}
\usepackage{amsmath,amssymb,mathtools}
\usepackage[scr=rsfs,cal=euler]{mathalfa}
\usepackage{xfrac}
\usepackage[unicode,hidelinks,bookmarks=false]{hyperref}
\newcommand{\ie}{i.e.\ }
\newcommand{\cf}{cf.\ }
\newcommand{\resp}{respectively\ }
\newcommand{\Subsection}{\subsection}
\providecommand{\nobreakdash}{\nobreak\hbox{-}}
\setlength{\emergencystretch}{2em}
\multlinegap0pt
\usepackage{bm}
\usepackage{bbm}
\usepackage{dsfont}
\usepackage{xspace}
\usepackage{cancel}
\usepackage{mathtools}
\usepackage{amsthm}
\makeatletter
\@ifundefined{theorem}{\newtheorem{theorem}{Theorem}}{}
\@ifundefined{lemma}{\newtheorem{lemma}[theorem]{Lemma}}{}
\@ifundefined{proposition}{\newtheorem{proposition}[theorem]{Proposition}}{}
\makeatother
\title[Universal Shuffle Asymptotics, Part II: Non-Gaussian Limits ]{Universal Shuffle Asymptotics, Part II: Non-Gaussian Limits for Shuffle Privacy -- Poisson, Skellam, and Compound-Poisson Regimes}
\author{Alex Shvets}
\date{Source: arXiv:2603.10073v1 (CC BY 4.0)}
\begin{document}
\maketitle

\noindent\emph{Source and licence.} Universal Shuffle Asymptotics, Part II: Non-Gaussian Limits for Shuffle Privacy -- Poisson, Skellam, and Compound-Poisson Regimes. Version arXiv:2603.10073v1 (CC BY 4.0), distributed under the Creative Commons Attribution 4.0 International licence, as recorded in the arXiv licence metadata for this exact version and shown on the abstract page https://arxiv.org/abs/2603.10073v1. Full rights notice: licenses/arxiv-2603.10073-CC-BY-4.0.txt.\par\smallskip

\noindent\textit{Editorial scope.} Counted unit: the Lemma whose source label is \texttt{lem:interior-conditional-smoothing} in the article (source0.tex lines 2301--2329, source label \texttt{lem:interior-conditional-smoothing}), delivered together with the article's own complete original proof of it (source0.tex lines 2331--2466, one \texttt{proof} environment of 136 lines). No mathematics is written, repaired or re-derived: statement and proof are the article's own text line for line. Required context retained uncounted: lem:tv-tensorization (lines 205--216); prop:hybrid-gaussian-compound-poisson (lines 1205--1257). Every definition, display and label of those lines is retained unchanged. Omitted from this package and not redistributed: 1--204, 217--1204, 1258--2300, 2330--2330, 2467--2585. Nothing inside a retained excerpt is omitted or altered: every retained body is delivered exactly as the article's lines are written, so no omission receipt of any kind accompanies this package. Citations: the retained excerpts carry no citation command at all, so no bibliography entry of the article is transcribed and no reference is redistributed. Reference adaptation: none was needed and none was made; every reference inside the delivered unit resolves inside it, so no unresolved reference or citation is produced. Printed slips kept literal: no printed word, symbol, hypothesis or number of the counted statement or of its proof was altered. Layout and class: the article's own class is \texttt{article} with packages including amsmath, amsthm, amssymb, hyperref, geometry, tcolorbox; the wrapper is the lane's pinned \texttt{amsart} editorial wrapper at 10pt on A4, which restates the theorem-like environments the retained text needs and adds the article's own macro definitions that the retained text needs (none), with extra wrapper packages (bm, bbm, dsfont, xspace, cancel, mathtools). The delivered numbering is the unit's own. Assets: no retained span contains a figure environment, a graphics-inclusion command, a PDF-page inclusion, a table or a TikZ picture; no image file is copied into this package, the package declares no asset dependency and no image file is redistributed with it. Licence: version arXiv:2603.10073v1 of this article is distributed under the Creative Commons Attribution 4.0 International licence, as recorded in the arXiv licence metadata for this exact version and shown on the abstract page https://arxiv.org/abs/2603.10073v1. A full rights notice is delivered at licenses/arxiv-2603.10073-CC-BY-4.0.txt. Characters: every retained excerpt is pure ASCII, so the delivered engine, XeLaTeX with its default Latin Modern text font, reports no missing character.
\begin{lemma}[Tensorization / union bound for independent pairs]
\label{lem:tv-tensorization}
Let $(X_1, Y_1)$ and $(X_2, Y_2)$ be couplings such that $X_1 \perp X_2$, $Y_1 \perp Y_2$, and the pair couplings
are independent across indices. Then
\[
\mathrm{TV}(\mathcal{L}(X_1, X_2), \mathcal{L}(Y_1, Y_2)) \leq P(X_1 \neq Y_1) + P(X_2 \neq Y_2).
\]
In particular, if $\mu_i, \nu_i$ are measures on $\mathcal{X}_i$ and $\mu = \mu_1 \otimes \mu_2$, $\nu = \nu_1 \otimes \nu_2$, then
\[
\mathrm{TV}(\mu, \nu) \leq \mathrm{TV}(\mu_1, \nu_1) + \mathrm{TV}(\mu_2, \nu_2).
\]
\end{lemma}

\begin{proposition}[Hybrid Gaussian / compound-Poisson limit]
\label{prop:hybrid-gaussian-compound-poisson}
Assume the two-dominant sparse-error critical regime above. Also assume
$D_0\cap D_1=\varnothing$. Finally suppose $\pi_n\to\pi\in[0,1]$. Define the Gaussian
covariance operator
\[
\Sigma:=(1-\pi)p_0(1-p_0)\,g_0g_0^{\top}+\pi p_1(1-p_1)\,g_1g_1^{\top}
\qquad \text{on } M,
\]
and, for each $y\notin D_b$, the jump vector
\[
j_{b,y}:=\Pi_J(e_y-\mu_b)\in M^\perp.
\]
Let $\nu$ be the finite measure on $M^\perp$ given by
\[
\nu
:=\sum_{y\notin D_0}(1-\pi)\alpha_0(y)\,\delta_{j_{0,y}}
 +\sum_{y\notin D_1}\pi\alpha_1(y)\,\delta_{j_{1,y}}.
\]
Let $G\sim N(0,\Sigma)$ and let $J$ be an independent compound-Poisson random vector with L\'evy measure
$\nu$, equivalently
\[
J\overset{d}=\sum_{y\notin D_0}U_y j_{0,y}+\sum_{y\notin D_1}V_y j_{1,y},
\]
where the coordinates are independent and
\[
U_y\sim \mathrm{Poi}((1-\pi)\alpha_0(y)),
\qquad
V_y\sim \mathrm{Poi}(\pi\alpha_1(y)).
\]
Set
\[
\Delta:=\Pi_J(\mu_1-\mu_0)\in M^\perp.
\]
Let $P_n$ be the law of $S_n$ under $T_{n,k_n}$, and let $Q_n$ be the law of the same statistic $S_n$
under $T_{n,k_n+1}$. Then
\[
P_n\Longrightarrow P_\infty:=\mathcal L(G,J),
\qquad
Q_n\Longrightarrow Q_\infty:=\mathcal L(G,J+\Delta).
\]
Equivalently, for $u\in M$ and $v\in M^\perp$,
\[
\mathbb E\exp\bigl(i\langle u,G\rangle+i\langle v,J\rangle\bigr)
=\exp\!\left(-\frac12\langle u,\Sigma u\rangle
+\int_{M^\perp}(e^{i\langle v,z\rangle}-1)\,\nu(dz)\right),
\]
whereas under the neighboring alternative the limiting characteristic function is multiplied by the shift factor
$e^{i\langle v,\Delta\rangle}$.

In particular, the limiting normalized experiment factors into a Gaussian component on the dominant block and an
independent compound-Poisson component on the rare-jump block.
\end{proposition}

\begin{lemma}[Interior conditional smoothing for the dominant block]
\label{lem:interior-conditional-smoothing}
Assume the setting of Proposition~\ref{prop:hybrid-gaussian-compound-poisson} and, in addition, $\pi_n\to\pi\in(0,1)$. Define
\[
\theta_{b,n}:=\frac{W_b^{(n)}(y_{ba})}{W_b^{(n)}(y_{ba})+W_b^{(n)}(y_{bb})},
\qquad b\in\{0,1\}.
\]
For each $z\in M^\perp$ in the support of $Z_n$, let
\[
T_{b,n}(z):=N_{n,k_n}(D_b),
\qquad b\in\{0,1\},
\]
noting that $T_{b,n}(z)$ is $\sigma(Z_n)$-measurable because $\ker \Pi_J=M=\mathrm{span}\{g_0,g_1\}$ and adding an element of $M$ changes only the within-pair differences on $D_0,D_1$, not the pair totals. Let $R_{n,z}$ be the law of
\[
\Gamma_{n,z}:=\Psi_{n,z}(X_{0,n}^\circ,X_{1,n}^\circ),
\]
where $X_{0,n}^\circ\sim \mathrm{Bin}(T_{0,n}(z),\theta_{0,n})$ and $X_{1,n}^\circ\sim \mathrm{Bin}(T_{1,n}(z),\theta_{1,n})$ are independent and
\[
\Psi_{n,z}(x_0,x_1):=n^{-1/2}\Bigl((x_0-p_0T_{0,n}(z))g_0+(x_1-p_1T_{1,n}(z))g_1\Bigr).
\]
Then there exists a finite constant $C_{\mathrm{int}}=C_{\mathrm{int}}(p_0,p_1,\pi,\{\alpha_b(y)\})$ such that, for all sufficiently large $n$,
\begin{equation}
\int \mathrm{TV}\bigl(Q_{n,z}^{G\mid Z},R_{n,z}\bigr)\,Q_n^J(dz)
+\int \mathrm{TV}\bigl(P_{n,z}^{G\mid Z},R_{n,z}\bigr)\,P_n^J(dz)
\le \frac{C_{\mathrm{int}}}{\sqrt n},
\label{eq:B3-main}
\end{equation}
where $P_{n,z}^{G\mid Z}$ and $Q_{n,z}^{G\mid Z}$ are regular conditional laws of $G_n$ given $Z_n=z$ under $T_{n,k_n}$ and $T_{n,k_n+1}$, respectively.
\end{lemma}

\begin{proof}
Fix one of the two hypotheses and let $m_{0,n}^\star,m_{1,n}^\star$ denote the corresponding group sizes:
\[
(m_{0,n}^\star,m_{1,n}^\star)=
\begin{cases}
(m_{0,n},m_{1,n}), & \text{under }T_{n,k_n},\\
(m_{0,n}-1,m_{1,n}+1), & \text{under }T_{n,k_n+1}.
\end{cases}
\]
Let $L_{0,n}$ be the number of $0$-input users whose output falls in $\mathcal Y\setminus D_0$, and let $L_{1,n}$ be the number of $1$-input users whose output falls in $\mathcal Y\setminus D_1$. Let $A_n$ be the number of $1$-input users whose output falls in $D_0$, and let $B_n$ be the number of $0$-input users whose output falls in $D_1$. Then $A_n\le L_{1,n}$ and $B_n\le L_{0,n}$. Conditioning on the full rare configuration (equivalently: on $Z_n=z$ together with the latent assignment of rare outputs to the two input groups), the counts
\[
X_{0,n}:=N_{n,k_n}(y_{0a}),
\qquad
X_{1,n}:=N_{n,k_n}(y_{1a})
\]
are independent and satisfy
\begin{align*}
X_{0,n}
&\overset d= \mathrm{Bin}(m_{0,n}^\star-L_{0,n},\theta_{0,n})+\mathrm{Bin}(A_n,\vartheta_{10,n}),\\
X_{1,n}
&\overset d= \mathrm{Bin}(m_{1,n}^\star-L_{1,n},\theta_{1,n})+\mathrm{Bin}(B_n,\vartheta_{01,n}),
\end{align*}
where
\[
\vartheta_{10,n}:=\frac{W_1^{(n)}(y_{0a})}{W_1^{(n)}(y_{0a})+W_1^{(n)}(y_{0b})},
\qquad
\vartheta_{01,n}:=\frac{W_0^{(n)}(y_{1a})}{W_0^{(n)}(y_{1a})+W_0^{(n)}(y_{1b})},
\]
with an arbitrary value when the denominator is zero (then the corresponding count $A_n$ or $B_n$ is automatically zero). Moreover,
\[
T_{0,n}(z)=(m_{0,n}^\star-L_{0,n})+A_n,
\qquad
T_{1,n}(z)=(m_{1,n}^\star-L_{1,n})+B_n.
\]
Hence $R_{n,z}$ is exactly the law obtained by replacing the cross terms $\mathrm{Bin}(A_n,\vartheta_{10,n})$ and $\mathrm{Bin}(B_n,\vartheta_{01,n})$ by the native terms $\mathrm{Bin}(A_n,\theta_{0,n})$ and $\mathrm{Bin}(B_n,\theta_{1,n})$.

Write
\[
S_{0,n}\sim \mathrm{Bin}(m_{0,n}^\star-L_{0,n},\theta_{0,n}),
\qquad
V_{0,n}\sim \mathrm{Bin}(A_n,\vartheta_{10,n}),
\qquad
V_{0,n}^\circ\sim \mathrm{Bin}(A_n,\theta_{0,n}),
\]
independently, and analogously define $S_{1,n},V_{1,n},V_{1,n}^\circ$ on the $D_1$-pair. Since
\[
S_{0,n}+V_{0,n}^\circ\sim \mathrm{Bin}(T_{0,n}(z),\theta_{0,n}),
\qquad
S_{1,n}+V_{1,n}^\circ\sim \mathrm{Bin}(T_{1,n}(z),\theta_{1,n}),
\]
it suffices, by contraction under the measurable map $\Psi_{n,z}$ and by Lemma~\ref{lem:tv-tensorization}, to control the two one-dimensional distances
\[
\mathrm{TV}\bigl(\mathcal L(S_{0,n}+V_{0,n}),\mathcal L(S_{0,n}+V_{0,n}^\circ)\bigr),
\qquad
\mathrm{TV}\bigl(\mathcal L(S_{1,n}+V_{1,n}),\mathcal L(S_{1,n}+V_{1,n}^\circ)\bigr).
\]
For the first term, couple $V_{0,n}$ and $V_{0,n}^\circ$ as sums of $A_n$ coupled Bernoulli pairs, one pair for each cross message, so that
\[
E\bigl|V_{0,n}-V_{0,n}^\circ\mid Z_n,\text{rare configuration}\bigr|
\le A_n\,|\vartheta_{10,n}-\theta_{0,n}|
\le A_n.
\]
Conditional on the coupled values $(V_{0,n},V_{0,n}^\circ)=(v,w)$, repeated one-step shifts give
\[
\mathrm{TV}\bigl(\mathcal L(S_{0,n}+v),\mathcal L(S_{0,n}+w)\bigr)
\le |v-w|\,\sup_k P(S_{0,n}=k).
\]
Taking conditional expectation yields
\[
\mathrm{TV}\bigl(\mathcal L(S_{0,n}+V_{0,n}),\mathcal L(S_{0,n}+V_{0,n}^\circ)\bigr)
\le A_n\,\sup_k P(S_{0,n}=k).
\]
The same argument on $D_1$ gives
\[
\mathrm{TV}\bigl(\mathcal L(S_{1,n}+V_{1,n}),\mathcal L(S_{1,n}+V_{1,n}^\circ)\bigr)
\le B_n\,\sup_k P(S_{1,n}=k).
\]
Now $\theta_{b,n}\to p_b\in(0,1)$, so for all large $n$,
\[
\theta_{b,n}(1-\theta_{b,n})\ge \frac12 p_b(1-p_b),
\qquad b\in\{0,1\},
\]
and the unimodal binomial mass bound gives
\[
\sup_k P(S_{b,n}=k)
\le \frac{c_b}{\sqrt{m_{b,n}^\star-L_{b,n}+1}},
\qquad
c_b:=\sqrt{\frac{2}{p_b(1-p_b)}}.
\]
Therefore, conditional on the rare configuration,
\begin{equation}
\mathrm{TV}\bigl(\mathcal L(G_n\mid Z_n,\text{rare configuration}),R_{n,Z_n}\bigr)
\le c_0\frac{A_n}{\sqrt{m_{0,n}^\star-L_{0,n}+1}}+c_1\frac{B_n}{\sqrt{m_{1,n}^\star-L_{1,n}+1}}.
\label{eq:B3-conditional-bound}
\end{equation}
Averaging over the latent rare configuration and then over $Z_n$ (tower property) gives the same bound for the conditional laws $P_{n,z}^{G\mid Z}$ or $Q_{n,z}^{G\mid Z}$.

It remains to integrate the right-hand side. Set
\[
\kappa:=\frac14\min\{\pi,1-\pi\}>0.
\]
Since $\pi_n\to\pi$, for all sufficiently large $n$ we have $m_{0,n}^\star\ge 3\kappa n$ and $m_{1,n}^\star\ge 3\kappa n$ under both hypotheses. Let $R_n:=L_{0,n}+L_{1,n}$. On the event $\{R_n\le \kappa n\}$,
\[
\frac{A_n}{\sqrt{m_{0,n}^\star-L_{0,n}+1}}+\frac{B_n}{\sqrt{m_{1,n}^\star-L_{1,n}+1}}
\le \frac{A_n+B_n}{\sqrt{2\kappa n}}
\le \frac{R_n}{\sqrt{2\kappa n}}.
\]
On the complementary event, the same quantity is bounded by $R_n$. Hence
\[
E\Bigl[\frac{A_n}{\sqrt{m_{0,n}^\star-L_{0,n}+1}}+\frac{B_n}{\sqrt{m_{1,n}^\star-L_{1,n}+1}}\Bigr]
\le \frac{E[R_n]}{\sqrt{2\kappa n}}+E\bigl[R_n\mathbf 1_{\{R_n>\kappa n\}}\bigr].
\]
Since $R_n$ is a Poisson-binomial sum of rare-output indicators,
\[
E[R_n]\le 2(\Lambda_0+\Lambda_1),
\qquad
E[R_n^2]\le 2(\Lambda_0+\Lambda_1)+4(\Lambda_0+\Lambda_1)^2
\]
for all large $n$, where
\[
\Lambda_b:=\sum_{y\notin D_b}\alpha_b(y), \qquad b\in\{0,1\}.
\]
By Cauchy--Schwarz and Markov,
\[
E\bigl[R_n\mathbf 1_{\{R_n>\kappa n\}}\bigr]
\le \frac{E[R_n^2]}{\kappa n}
=O(n^{-1}).
\]
Combining this with \eqref{eq:B3-conditional-bound} gives \eqref{eq:B3-main}. One may take, for example,
\[
C_{\mathrm{int}}
:=\Bigl(\sqrt{\frac{2}{p_0(1-p_0)}}+\sqrt{\frac{2}{p_1(1-p_1)}}\Bigr)
\Bigl(\frac{2(\Lambda_0+\Lambda_1)}{\sqrt{2\kappa}}+\frac{2(\Lambda_0+\Lambda_1)+4(\Lambda_0+\Lambda_1)^2}{\kappa}\Bigr).
\]
Since $\mathcal Y$ is finite, this depends only on $p_0,p_1,\pi$ and finitely many $\alpha_b(y)$; in particular it is controlled by $p_0,p_1,\pi$ and $\sup_{b,y}\alpha_b(y)$ once $\mathcal Y$ is fixed.
\end{proof}

\end{document}
